\documentclass[11pt, oneside]{article} 
\usepackage{geometry}
\geometry{letterpaper} 
\usepackage{graphicx}
	
\usepackage{amssymb}
\usepackage{amsmath}
\usepackage{parskip}
\usepackage{color}
\usepackage{hyperref}

\graphicspath{{/Users/telliott/Dropbox/Github-math/figures/}}
% \begin{center} \includegraphics [scale=0.4] {gauss3.png} \end{center}

\title{Menelaus's theorem}
\date{}

\begin{document}
\maketitle
\Large

%[my-super-duper-separator]

$\circ \ $ \ Medians divide the sides in half, and they cross at the centroid.

$\circ \ $ \ Altitudes meet the sides (or their extensions) at right angles, and they cross at the orthocenter.

$\circ \ $ \ The angle bisectors meet at the incenter, the center of the incircle.

Any two lines which are not identical or parallel cross at a point.  Here we are concerned to prove that the third median crosses at the same point --- the three lines are concurrent.

\label{sec:midpoint_theorem}

We saw the \hyperref[sec:midline_theorem]{\textbf{midline theorem}}  earlier:

$\bullet$ \ The midline of a triangle is parallel to the third side of that triangle and one-half its length.

\begin{center} \includegraphics [scale=0.18] {midline_thm.png} \end{center}

This theorem is actually a special case of the general theorems about similar triangles.  Just as for the similarity theorem, the converse is also true.  

\subsection*{application to Ceva's theorem}

\begin{center} \includegraphics [scale=0.3] {centroid_ratio.png} \end{center}

\emph{Proof}.

In $\triangle ABC$ draw the medians $AK$ and $BL$ and find the centroid $G$ where the two medians cross.

Draw $CG$ and extend it across the base at $M$ to meet $AP \parallel LGB$.

Because $LGB \parallel AP$ we have $\triangle CLG \sim \triangle CAP$ in the ratio 1:2.

It follows that $CG = 1/2 \cdot CP$.

But we are given $CK = 1/2 \cdot CB$.

It follows that $\triangle CKG \sim \triangle CBP$ in the ratio 1:2.

By the midline theorem, $AGK \parallel BP$.

Thus $AGBP$ is a parallelogram.

In a parallelogram the diagonals bisect each other so $AM = BM$.

$\square$

This proof is not quite as elegant as the one that comes from Menelaus's theorem.  But the parallelogram provides a very nice, simple approach to the division of the medians at the centroid $G$.  

You can think of the centroid as the center of gravity of $\triangle ABC$.  If it were constructed of a sheet of constant thickness and balanced on the head of a pin, the pin would be at the point $G$.

\subsection*{position of centroid}

The centroid divides the medians in the ratio 2:1.  

\emph{Proof}.

\begin{center} \includegraphics [scale=0.3] {centroid_ratio.png} \end{center}

Given medians and centroid for $\triangle ABC$.

Extend $GM$ to form $PM = GM$.  

Since the diagonals are bisected, $AGBP$ is a parallelogram.

Since $LGB \parallel AP$ and $CL = AL$, we have $\triangle CLG \sim \triangle CAP$ with ratio 1/2 

Thus $CG = GP$.

$GM$ is $1/2 \cdot GP$. so it is also $1/2 \cdot CG$.

$GM$ is 1/3 of the whole length of the median $CM$.

$\square$

\subsection*{area proof}

The medians are lines from each vertex to the midpoint of the opposing side.  

\begin{center} \includegraphics [scale=0.16] {centroid4.png} \end{center}

Ceva's theorem says that if the lines cross at one point then $a_1/a_2 \cdot b_1/b_2 \cdot c_1/c_2 = 1$.  Here we use the converse:  given that the ratios multiply to give 1, the points are concurrent.

For medians each term like $a_1/a_2$ is just 1, so that part is easy.

As we said, the point where the medians meet is called the centroid and is usually given the label $G$.

$\bullet \ $The medians divide any triangle into six small triangles, all with equal area.

\begin{center} \includegraphics [scale=0.16] {centroid4.png} \end{center}

\emph{Proof}.

Write $(\triangle ABC)$ for the area of the triangle.

$(\triangle GBX) = (\triangle GCX)$ by the \hyperref[sec:area_ratio_theorem]{\textbf{area-ratio theorem}}.  They share the same vertex and have equal bases.  Hence the two small triangles on any one side have equal area:  $x = x'$, $y = y'$, $z = z'$.

Also by the same theorem, $(\triangle ABX) = (\triangle ACX)$.  So by subtraction $(\triangle ABG) = (\triangle ACG)$:  $y + y' = z + z'$.  Thus $y = z$.  But this is true for any side.  Hence all six triangles are equal in area.

$\square$

\begin{center} \includegraphics [scale=0.16] {centroid4.png} \end{center}

Consider $\triangle GBX$ and $\triangle ABX$ on base $XGA$.  They have the same altitude, from the base to the vertex at $A$.  The ratio of areas is $1:3$.  So that must also be the ratio $XG:XA$.

The centroid lies one-third of the way up the median from the side.

$\square$

\subsection*{Altitudes}

We had a nice (simple) \hyperref[sec:Newton_altitude]{\textbf{proof}}  from Newton that the altitudes of a triangle are concurrent, earlier in the book.  

Here is another one, based on the converse of Ceva's theorem.

Draw the altitudes in a triangle.  We can show that they are concurrent at a point, called the orthocenter.

This forms three pairs of similar triangles.  Form the ratios of parts of sides to parts of altitudes.

\begin{center} \includegraphics [scale=0.3] {altitudes_Ceva.png} \end{center}
\[ \frac{a}{b'}  = \frac{e}{g} \ \ \ \ \ \ \ \ \  \frac{b}{c'} =  \frac{g}{j} \ \ \ \ \ \ \ \ \ \frac{c}{a'} = \frac{j}{e}  \]

Multiply together all the terms on the left and separately the terms on the right.  They must be equal, and we can see that the right-hand side cancels to give 1.  Thus
\[ \frac{a}{b'} \cdot \frac{b}{c'} \cdot \frac{c}{a'} = 1 \]
\[ \frac{a}{a'} \cdot \frac{b}{b'} \cdot \frac{c}{c'} = 1 \]

By the converse of Ceva's theorem, we have that the altitudes are concurrent.

$\square$

\subsection*{problem}

Use the angle bisector theorem and the converse of Ceva's theorem to show that bisectors of the angles in a triangle meet at a point.

That point is called the incenter.

\end{document}